\begin{figure}[t]
\centering
\setlength{\fboxsep}{8pt}
\fbox{\begin{minipage}[t][100pt][t]{0.445\linewidth}
\small
\textbf{One program, more executions}\par\medskip
\centering
$\boxed{H_0^{(1)}}\qquad\boxed{H_0^{(2)}}$\par\smallskip
Same code; independent errors\par\medskip
$p=0.5\quad\longrightarrow\quad1-(1-p)^2=0.75$\par\medskip
\raggedright
\textbf{Coverage increases by $25$\,\%.}\par
Program specialization remains zero.
\end{minipage}}\hfill
\fbox{\begin{minipage}[t][100pt][t]{0.445\linewidth}
\small
\textbf{A measured same-code reference}\par\medskip
\centering
MATH-500: nine identical \bare{} slots\par\smallskip
Three independent repeats per slot\par\medskip
$H_{\rm stable}=0\qquad\widehat H=\WpVthreeClonePlugin\,\%$\par\medskip
\raggedright
\textbf{The observed estimate is positive.}\par
Descriptive plugin on $383$ complete tasks.
\end{minipage}}
\caption{\textbf{More coverage does not identify more specialization.}
Left: a probability calculation for two independent executions of one
program. Right: the empirical clone arm produces a positive
repeat-averaged oracle estimate despite zero program specialization
(complete-task estimate; \S\ref{sec:f3-clones}).
The panels illustrate different quantities, not matched effect sizes.
Our control design compares eight generated programs plus \bare{} with
nine identical \bare{} slots, then separates repeat validation from
feature-based pre-execution selection.}
\label{fig:sampling-control}
\end{figure}
